% ATTEMPT_STATUS: unresolved
% TOP_PROBLEM: {"problemId": "problem.asymptotic-exponent-of-matrix-multiplication-omega-2", "canonicalTitle": "Matrix-Multiplication Exponent Equals Two", "releaseRank": 26, "primaryDomain": "theoretical_computer_science", "primaryDomainLabel": "Theoretical computer science", "displayStatus": "open", "releaseStatus": "open"}
% !TeX program = lualatex
\documentclass[11pt]{article}
\usepackage[margin=25mm]{geometry}
\usepackage{fontspec}
\setmainfont{Latin Modern Roman}
\usepackage{amsmath,amssymb,mathtools}
\usepackage{unicode-math}
\setmathfont{Latin Modern Math}
\usepackage{xurl,hyperref}
\hypersetup{colorlinks=true,urlcolor=blue}
\setcounter{secnumdepth}{0}
\setlength{\parindent}{0pt}
\setlength{\parskip}{5pt}
\setlength{\emergencystretch}{3em}
\begin{document}
\section*{\texorpdfstring{Matrix-Multiplication Exponent Equals Two}{Matrix-Multiplication Exponent Equals Two}}\label{26}
\subsection{Definitions and mathematical statement}
The rational bilinear rank $R_{{\mathbb{Q}}}(n)$ is the least $r$ admitting rational linear forms $\alpha_j,\beta_j$ and fixed rational matrices $C_j$ with
\[
 AB=\sum_{j=1}^r\alpha_j(A)\beta_j(B)C_j\qquad(A,B\in{\mathbb{Q}}^{n\times n}).
\]
Set $\omega_{{\mathbb{Q}}}=\inf\{\tau:R_{{\mathbb{Q}}}(n)=O(n^\tau)\}$. The target is $\omega_{{\mathbb{Q}}}=2$, or equivalently the existence, for every $\varepsilon>0$, of exact arithmetic circuits of size $O_\varepsilon(n^{2+\varepsilon})$. This counts arithmetic operations at unit cost. It neither bounds coefficient bit lengths nor asserts that the infimum is attained by an $O(n^2)$ algorithm.
\par\smallskip\textit{Formulation and concept references:} \href{https://theoryofcomputing.org/articles/gs005/gs005.pdf}{[S1]}.

\subsection{Short English statement}
Can exact rational matrix multiplication be performed with an operation count arbitrarily close to quadratic?
\subsection{Research attempt}

\paragraph{Scope and source check (27 September 2026).}
All coefficients below lie in $\mathbb Q$, and arithmetic operations have
unit cost. For every $\varepsilon>0$ the target permits a different
algorithm and a different implied constant. It does not demand one
quadratic algorithm, a practical implementation, or a bound on the bit
length of its constants. Bl\"aser [S1] gives the classical rank and
approximation framework. The August 2026 paper [S2] reports
$\omega<2.371177$ and describes a rigorous verification procedure;
its numerical certificates were not rerun here. That bound is above two.
The infinite-group construction [S3, Section 3] explicitly identifies
dimension deficits preventing its construction from reaching two.
Neither reference is used as an assertion that the target is settled.

The attempt is to turn the desired asymptotic statement into exact finite
certificates, then determine whether approximation can supply them.
The arguments reproduce elementary parts of the classical framework,
including the principle behind Bini's approximation method [E1]. They
are not claimed as new matrix multiplication algorithms. A small tensor
example will test the difference between an exact identity and a limit.

\paragraph{Rank bounds and the block construction.}
Abbreviate $R_{\mathbb Q}(n)$ to $R(n)$. The output polynomials
$p_{ij}=\sum_k a_{ik}b_{kj}$ are linearly independent: the monomial
$a_{i1}b_{1j}$ occurs in exactly one of them. In a rank-$r$ algorithm
all these polynomials lie in the span of the $r$ products
$\alpha_s(A)\beta_s(B)$. Hence
\begin{equation}\label{mm:elementary}
 n^2\le R(n)\le n^3.
\end{equation}
The upper bound uses the individual products $a_{ik}b_{kj}$.
Padding inputs with zero rows and columns proves $R(n)\le R(N)$ for
$n\le N$.

More significantly,
\begin{equation}\label{mm:submult}
 R(ab)\le R(a)R(b).
\end{equation}
Apply a rank-$R(a)$ identity to $a$ by $a$ block matrices with $b$ by
$b$ blocks. Each block product is evaluated with a rank-$R(b)$ identity.
The resulting scalar multiplications are products of a rational linear
form in the first input and one in the second. The outer identity remains
valid for matrix blocks: after expansion every term has one first-input
block followed by one second-input block, in that order. Coefficient
comparison, rather than an assumption that blocks commute, proves its
validity. This gives the stated product of ranks.

Fix $b\ge2$ and a rank-$r$ identity for size $b$. Recursion and padding
give, for $k=\lceil\log_b n\rceil$,
\[
 R(n)\le r^k\le r n^{\log_b r}.
\]
It follows that $\omega_{\mathbb Q}\le\log_b r$. Conversely, if
$R(n)\le Cn^\tau$ for all sufficiently large $n$, then
$\log_bR(b)\le\tau+\log_b C$ for sufficiently large $b$.
Taking $b$ to infinity and then the infimum over $\tau$ proves
\begin{equation}\label{mm:finite}
 \omega_{\mathbb Q}=\inf_{b\ge2}\log_b R(b).
\end{equation}
In particular, exponent two is equivalent to the following exact finite
certificate condition:
\[
 \forall\delta>0\quad\exists b\ge2\quad
 \exists\text{ a rational rank identity with }r<b^{2+\delta}.
\]
Strict inequality causes no difficulty: use half the requested exponent
slack before absorbing a constant. There is no asserted bound on the
size of the certificate or on how rapidly its block size grows as
$\delta$ decreases.

\paragraph{Accounting for the linear operations.}
A rank count ignores additions, but the fixed-block recursion restores
their cost. All linear combinations in a chosen size-$b$ identity have
a finite cost $c_b$. For inputs of size $b^k$, their total cost at the
top level is at most $c_b b^{2(k-1)}$. Thus the operation count satisfies
\[
 T(b^k)\le rT(b^{k-1})+c_b b^{2(k-1)}.
\]
If $r>b^2$, summing the geometric series gives $T(n)=O_b(n^{\log_b r})$.
If $r=b^2$, the same calculation gives $T(n)=O_b(n^2\log n)$.
Both bounds imply the required $n^{2+\varepsilon}$ cost whenever a
suitable certificate from \eqref{mm:finite} has been chosen with slack.
The exceptional equality case must not be silently advertised as
quadratic total cost. No uniform bound on $c_b$ is needed after $b$ is
fixed for the chosen $\varepsilon$.

The distinction between a fixed saving and exponent two is now precise.
A certificate $r\le b^3/c$ yields the bound
$3-\log_b c$ on the exponent. If $1<c<b$, this particular bound exceeds
two. It does not show that the true exponent exceeds two. Likewise a
single bound $r\le Cb^2$, with $C>1$, only gives
$2+\log_b C$. Such bounds for unbounded $b$ with the \emph{same} $C$
would suffice, as would $r\le b^2\exp(o(\log b))$ along an unbounded
sequence. A table of a few improved ranks does not provide either family.

\paragraph{A rational degeneration can be converted exactly.}
Let $M_b$ denote the bilinear tensor for size-$b$ multiplication.
Suppose an explicit identity over $\mathbb Q[t]$ has the form
\begin{equation}\label{mm:degeneration}
 P(t)=\sum_{s=1}^r u_s(t)\otimes v_s(t)\otimes w_s(t)
     =t^qM_b+\sum_{j=q+1}^{D}t^j E_j,
 \qquad 0\le q\le D.
\end{equation}
All coefficients below degree $q$ vanish. These are formal coefficient
identities, not estimates obtained by taking $t$ numerically small.
Tensoring $k$ copies gives a polynomial of degree at most $kD$ whose
coefficient of $t^{kq}$ is $M_b^{\otimes k}=M_{b^k}$, after reindexing.

Choose $kD+1$ distinct rational numbers $a_0,\ldots,a_{kD}$.
Lagrange interpolation provides rational weights $\lambda_i$ such that
the coefficient of $t^{kq}$ in every polynomial of degree at most $kD$
equals $\sum_i\lambda_i F(a_i)$. This holds coordinatewise for tensors.
At each $a_i$, the tensor $P(a_i)^{\otimes k}$ has rank at most $r^k$.
Consequently
\begin{equation}\label{mm:extract}
 R(b^k)\le(kD+1)r^k,
 \qquad
 \omega_{\mathbb Q}\le\log_b r.
\end{equation}
For the second assertion, apply \eqref{mm:finite} to block size $b^k$
and let $k$ grow: $\log_{b^k}(kD+1)$ tends to zero. Alternatively,
padding gives a logarithmic factor in the rank bound, absorbed by every
positive exponent slack.

Here $b,r,D,q$ are fixed before the tensor power grows. A large degree
$D$ can make the practical certificate enormous without changing this
limit. For a family with varying $b$, first choose the certificate and
then take its tensor powers. Interchanging these choices without
controlling the new overhead would be unjustified. Laurent-polynomial
identities can be treated by clearing powers of $t$ first; an unspecified
analytic limit or a floating-point decomposition supplies neither the
coefficient identities nor their rationality.

There is also a distinction in the order of operations. Extracting an
exact rank bound $R(b)\le(D+1)r$ first and tensoring it produces
$((D+1)r)^k$. Extracting after tensoring, as above, costs only
$(kD+1)r^k$. The former loses a fixed amount in the exponent; the latter
does not. This calculation explains why a degeneration can matter even
when its first exact specialization is worse than a direct algorithm.

\paragraph{What algebraic coefficients cost over the rationals.}
The field requirement can also be audited quantitatively. Let
$E/\mathbb Q$ be a finite extension of degree $h$, and suppose a
rational bilinear map $T$ has a rank-$r$ decomposition over $E$.
Choose a rational basis $e_1,\ldots,e_h$ of $E$ and a rational-linear
map $\lambda:E\to\mathbb Q$ with $\lambda(1)=1$.
On rational inputs expand
$\alpha_s=\sum_u e_u\alpha_{su}$ and
$\beta_s=\sum_v e_v\beta_{sv}$, where the new forms are rational.
Applying $\lambda$ coordinatewise to the output gives the exact identity
\[
 T(A,B)=\sum_{s,u,v}\alpha_{su}(A)\beta_{sv}(B)
                          \lambda(e_ue_v C_s).
\]
Thus $R_{\mathbb Q}(T)\le h^2R_E(T)$. In particular, tensor powers
of one fixed size-$b$ algorithm can all be descended through the
\emph{same} extension, yielding $R_{\mathbb Q}(b^k)\le h^2r^k$.
The factor is $h^2$, not $h^{2k}$. The same argument after interpolation
in \eqref{mm:extract} gives $h^2(kD+1)r^k$ for a fixed degeneration
over $E$. Neither finite extension changes the resulting exponent bound.

This observation does not authorize treating a numerical complex
decomposition as an exact rational algorithm. Its coefficients must be
specified exactly, its polynomial identities checked, and the finite
extension hypothesis established before the displayed descent applies.
For any one certificate consisting of algebraic numbers, their
compositum is finite, so it qualifies. The degree may depend on the
certificate chosen for a particular exponent slack. The proof still
fixes that certificate before increasing its tensor power. No claim
about coefficient bit complexity follows from this operation count.

\paragraph{A small tensor exposes the limit trap.}
Let $e_0,e_1$ be a basis of $\mathbb Q^2$ and put
\[
 W=e_1\otimes e_0\otimes e_0+
   e_0\otimes e_1\otimes e_0+
   e_0\otimes e_0\otimes e_1.
\]
The polynomial
\[
 (e_0+te_1)^{\otimes3}-e_0^{\otimes3}
       =tW+t^2V+t^3e_1^{\otimes3}
\]
is a sum of two decomposable tensors at every rational $t$.
Nevertheless $R(W)=3$. For a direct proof, its first-factor slice space
is spanned by
\[
 A=\begin{pmatrix}0&1\\1&0\end{pmatrix},\qquad
 B=\begin{pmatrix}1&0\\0&0\end{pmatrix}.
\]
Since $\det(aA+bB)=-a^2$, all rank-one matrices in this space lie in
the single line spanned by $B$. A rank-two decomposition of $W$ would
force this two-dimensional slice space to be spanned by two rank-one
matrices: the two first-factor vectors must be independent, or the
slice space would have dimension at most one. This is impossible.
The displayed three-term expression gives the reverse bound.

Thus the limiting tensor need not have the ranks of the nearby tensors.
Coefficient extraction nevertheless proves
$R(W^{\otimes k})\le(3k+1)2^k$. This example is not a matrix
multiplication construction and does not establish the desired ranks
for $M_b$. Its purpose is to verify exactly what the approximation
argument supplies, and why substituting $t=0$ into a normalized
expression is not an exact rank-two algorithm.

\paragraph{Certificates, verification, and the remaining construction.}
For a proposed exact algorithm, expand its rational coefficient arrays
$\alpha_{s,ik},\beta_{s,\ell j},C_{s,pq}$. It suffices to check all
$b^6$ equations
\[
 \sum_{s=1}^r\alpha_{s,ik}\beta_{s,\ell j}C_{s,pq}
       =\delta_{ip}\delta_{jq}\delta_{k\ell}.
\]
These are equivalent to testing every pair of matrix units in the two
inputs. The local check verified all 64 scalar equations for the usual
seven-product size-two identity, including the signs in the four output
combinations. For the tensor $W$, rational interpolation was checked on
the first three tensor powers, against explicit coefficient expansion.
These finite checks support the stated identities; they are not a
search result for near-quadratic algorithms.

Even monotonicity, submultiplicativity, and an infimum exponent of two
do not formally force an attained quadratic bound. As an abstract test,
the integer sequence
\[
 a(n)=\left\lceil n^2\exp(\sqrt{\log n})\right\rceil
\]
is monotone and submultiplicative, has infimum exponent two, and is not
$O(n^2)$. Submultiplicativity follows from
$\sqrt{x+y}\le\sqrt x+\sqrt y$ and
$\lceil uv\rceil\le\lceil u\rceil\lceil v\rceil$ for positive $u,v$.
This sequence is not asserted to be a possible matrix multiplication
rank function; it only disproves that proposed deduction from those
abstract properties alone.

The first missing step is an unbounded family of rational exact
identities, or explicit degenerations \eqref{mm:degeneration}, with
$\log_b r$ approaching two. Neither the block lemma nor interpolation
constructs this family. Three concrete next tasks are to search for
small rational degenerations with exact residual checking; to prove a
structural rank saving that improves with block size; and to track the
actual degrees and coefficient fields in any proposed infinite-group
construction before applying \eqref{mm:extract}.

\textbf{UNRESOLVED.} The attempt gives finite sufficient certificates,
a complete rational approximation-to-exact argument, and countertests
to several stronger conclusions. It supplies no new exponent bound.
Written by GPT-6 Astra Ultra with source inspection and exact finite
algebra checks; no proof-assistant verification is claimed.

\subsection{Sources}
\begin{itemize}
\item[S1]Markus Bläser, Fast Matrix Multiplication, Theory of Computing Graduate Surveys 5 (2013).
\url{https://theoryofcomputing.org/articles/gs005/gs005.pdf}.
\textit{Formulation source.}
\item[S2]Improving the matrix multiplication exponent with modern optimization and AlphaEvolve.
\url{https://arxiv.org/abs/2608.16884}.
\textit{Further reference; not a proof endorsement.}
\item[S3]Finite Matrix Multiplication Algorithms from Infinite Groups.
\url{https://drops.dagstuhl.de/storage/00lipics/lipics-vol325-itcs2025/html/LIPIcs.ITCS.2025.18/LIPIcs.ITCS.2025.18.html}.
\textit{Further reference; not a proof endorsement.}
\item[E1]Dario Bini, \emph{Relations between exact and approximate bilinear algorithms. Applications}, Calcolo 17 (1980), 87--97.
\url{https://doi.org/10.1007/BF02575865}.
Also Dario Bini, Grazia Lotti, and Francesco Romani, \emph{Approximate Solutions for the Bilinear Form Computational Problem}, SIAM Journal on Computing 9 (1980), 692--697.
\url{https://epubs.siam.org/doi/10.1137/0209053}.
\textit{Historical primary references for approximation methods; bibliographic record and second abstract consulted 27 September 2026. The rational interpolation argument used here is given in full.}
\end{itemize}
\end{document}
